\documentclass[11pt,a4paper]{article}
\usepackage[margin=1in]{geometry}
\usepackage{amsmath,amssymb}
\usepackage{booktabs}
\usepackage{array}
\usepackage[T1]{fontenc}
\usepackage{lmodern}
\usepackage{microtype}
\usepackage{textcomp}
\usepackage{xcolor}
\definecolor{seclink}{HTML}{1F4E8C}
\usepackage[colorlinks=true,linkcolor=seclink,citecolor=seclink,urlcolor=blue]{hyperref}
\usepackage{parskip}
\usepackage[most]{tcolorbox}

\definecolor{ink}{HTML}{1A1A1A}
\definecolor{grey}{HTML}{5F5B54}
\definecolor{yellowbg}{HTML}{FBF3D5}
\definecolor{pinkbg}{HTML}{FADDE1}

\tcbset{plainbg/.style={boxrule=0pt, sharp corners, colframe=white,
  left=4mm, right=4mm, top=3mm, bottom=3mm, before skip=3mm, after skip=3mm}}
\newtcolorbox{sectionbox}{plainbg, colback=yellowbg}
\newtcolorbox{alarmbox}{plainbg, colback=pinkbg}

\title{\vspace{-2em}The Invisibles}
\author{\normalsize The largest companies nobody has heard of, and why that is structural}
\date{\normalsize Claude, 6 September 2026}

\begin{document}
\maketitle
\vspace{-1em}

\begin{sectionbox}
Robin asked which multinationals are enormous and unheard of, and offered Unilever
as the reference. Unilever is the mild case. It is a name most people half
recognise attached to products they use every day without connecting the two.
There are firms considerably further along that axis, and the interesting part is
that their obscurity is not an accident of marketing. It follows from where they
sit in the composition.
\end{sectionbox}

\section{Invisible by position}

These never sell to a person, so no person has any reason to learn the name.

\begin{center}
\begin{tabular}{@{}l@{\hspace{5mm}}l@{\hspace{5mm}}>{\raggedright\arraybackslash}p{0.50\textwidth}@{}}
\toprule
\textbf{Firm} & \textbf{Revenue} & \textbf{What it is} \\
\midrule
Vitol      & \$250--400bn & The largest independent oil trader. Private. Revenue swings
                            enormously with the oil price \\
Trafigura  & \$250bn+     & The same trade, plus metals \\
McKesson   & $\approx$\$310bn & United States pharmaceutical distribution. Among the largest
                            American companies by revenue \\
Cencora    & $\approx$\$290bn & The second of three. With Cardinal Health these firms move
                            close to 90\% of American prescription drugs \\
Cargill    & \$160--180bn & The largest private company in the United States. Grain, meat
                            and salt. Held by one family \\
Hon Hai    & $\approx$\$215bn & Foxconn. It assembles the iPhone \\
MSC        & undisclosed  & The largest container shipping line in the world, ahead of
                            Maersk. Private, held by the Aponte family \\
Compass    & $\approx$\pounds31bn & Feeds people in offices, hospitals, schools, stadiums and
                            oil rigs across some thirty countries \\
Munich Re  & $\approx$\texteuro60bn & Reinsurance. It insures the insurers, and therefore prices
                            the tail risk of civilisation, including the climate \\
\bottomrule
\end{tabular}
\end{center}

\section{The Unilever pattern}

Here the products are famous and the parent is not.

\begin{center}
\begin{tabular}{@{}l@{\hspace{5mm}}>{\raggedright\arraybackslash}p{0.63\textwidth}@{}}
\toprule
\textbf{Firm} & \textbf{What it holds} \\
\midrule
EssilorLuxottica & Ray-Ban, Oakley and Persol, plus the eyewear licences for Chanel,
                   Prada, Armani, Versace and Burberry. It also owns Sunglass Hut,
                   LensCrafters, Pearle Vision and Vision Express, and an eyecare
                   insurer \\[1mm]
Givaudan         & With dsm-firmenich, IFF and Symrise, four firms substantially
                   determine what processed food tastes like and what nearly
                   everything smells like \\[1mm]
JAB Holding      & The Reimann family. Panera, Krispy Kreme, Pret A Manger, Peet's,
                   and a large position in Keurig Dr Pepper \\[1mm]
Reckitt          & The medicine cabinet. Dettol, Durex, Nurofen, Strepsils, Finish,
                   Air Wick \\[1mm]
Henkel           & Persil, Schwarzkopf and Loctite \\[1mm]
Richemont        & Cartier, Van Cleef \& Arpels, Montblanc and IWC \\[1mm]
Wilmar           & A very large share of the world's palm and cooking oil \\
\bottomrule
\end{tabular}
\end{center}

EssilorLuxottica is the strongest single answer to the question. It designs the
frames, manufactures them, licenses the famous names onto them, sells them in its
own retail chains, and insures the purchase. That is a near complete vertical hold
on how human beings correct their vision, and the name means nothing to almost
anybody who wears the product.

\section{Why this happens}

Public recognition is bought with advertising, and advertising is only worth buying
when a person chooses you. None of the firms above are chosen by a person. So
public awareness of the economy runs roughly inversely to economic weight, and the
inversion is not a quirk. It is a consequence of where the last transaction sits.

\section{The container reading}

\begin{sectionbox}
The economy is a composition of interfaces, and a consumer stands at the outer
boundary of it. He can query exactly one interface, the last one in the chain.
Everything upstream has been composed away.

That is not a defect of the composition. It is what composition is for. A morphism
gives a forward mapping on prompts and a backward mapping on replies, and it tells
you nothing whatsoever about what stands behind it. The interface presented is
\emph{Ray-Ban}. Nothing in the type says EssilorLuxottica, and nothing needs to.
\end{sectionbox}

Which makes this the same observation as the other two notes, arriving from a third
direction. \texttt{terror.txt} concerns a centre that cannot see what its own
commissariats know. \texttt{anarchy.txt} concerns a federation that accidentally
concentrates knowledge in whoever gets re-mandated. This concerns an economy in
which the largest structures are unobservable from the only position most people
can occupy.

In each case the question is the same one. Who can see what, and what does the
composition hide from whom. It is a pleasing feature of your algebra that it treats
concealment as a structural property rather than as a moral failing, because in two
of these three cases nobody is doing anything wrong and the opacity is there
regardless.

\begin{alarmbox}
\textbf{On the numbers.} Treat every figure as approximate. Commodity trading
revenue in particular moves by hundreds of billions with the oil price, private
firms disclose selectively or not at all, and my information has a horizon around
May 2026. The ranking is robust. The digits are not.
\end{alarmbox}

\end{document}
